\documentclass[14pt]{extarticle}
\usepackage{fontspec}
\usepackage{polyglossia}
\usepackage{amsmath}
\usepackage{amssymb}
\usepackage{xcolor}
\usepackage{fancyhdr}
\usepackage{graphicx}
\usepackage{listings}
\usepackage[most]{tcolorbox}
% \usepackage{setspace}
\usepackage{pifont}
\usepackage{enumitem}
\usepackage{titlesec}
\usepackage[bottom]{footmisc}
\usepackage{titling}
\usepackage{minted}
\usepackage{etoolbox}
\usepackage{array}
\usepackage{longtable} % page-breakable tables (must load before polyglossia/bidi RTL setup)
\usepackage{extsizes}
\usepackage{float}
\usepackage{booktabs}
\usepackage{hyperref}
\usepackage[onehalfspacing]{setspace}
\newif\ifwhiteblack
\ifdefined\setwhiteblack
  \setwhiteblack
\fi

\ifwhiteblack
  \usepackage[a4paper, landscape, top=1.25cm, bottom=1.25cm, left=1.25cm, right=1.25cm, headheight=1.5cm, headsep=0.5cm]{geometry}
\else
  \usepackage[a4paper, top=2.7cm, bottom=2.7cm, left=2.5cm, right=2.5cm, headheight=1.5cm, headsep=0.5cm]{geometry}
\fi
\usepackage{tikz}
\usepackage{datetime2}

\usetikzlibrary{automata, positioning, arrows.meta, arrows}
\providecommand{\LTcaptype}{table}

\newfontfamily\emoji{Segoe UI Emoji}

\pagestyle{fancy}

\setmainlanguage[numerals=western]{arabic}
\setotherlanguage{english}
\newfontfamily\arabicfont[Script=Arabic]{Calibri}
\newfontfamily\arabicfonttt[Script=Arabic]{Courier New}



\ifwhiteblack
\lstset{
  language=[Sharp]C,
  numbers=left,
  stepnumber=1,
  numbersep=8pt,
  frame=single,
  basicstyle=\ttfamily\small,
  keywordstyle=\bfseries,
  stringstyle=\ttfamily,
  commentstyle=\itshape
}
\else
\lstset{
  language=[Sharp]C,
  numbers=left,
  stepnumber=1,
  numbersep=8pt,
  frame=single,
  basicstyle=\ttfamily\small,
  keywordstyle=\color{blue},
  stringstyle=\color{red},
  commentstyle=\color{green!50!black}
}
\fi

\newif\ifdetailed
\ifdefined\setdetailed
  \setdetailed
\fi

\newif\ifwithsols
\ifdefined\setwithsols
  \setwithsols
\fi

\ifwhiteblack
  \AtBeginDocument{
    \LARGE
    \usemintedstyle{bw}
  }
\fi

% Default minted look: compact code that wraps long lines.
% Not applied in black-and-white landscape mode, which already uses \LARGE.
\ifwhiteblack\else
  \setminted{fontsize=\small, breaklines=true}
\fi

% unified tcolorboxes for articles
\ifwhiteblack
\tcbset{colback=white, colframe=black, fonttitle=\bfseries, boxrule=1pt}
\newtcolorbox{boxDef}[1][]{title={{\emoji📘} تعريف\ifx\\#1\\\else ~#1\fi :}}
\newtcolorbox{boxExercise}[1][]{title={{\emoji🧩} تمرين\ifx\\#1\\\else ~#1\fi :}}
\newtcolorbox{boxExample}[1][]{title={{\emoji📝} مثال\ifx\\#1\\\else ~#1\fi :}}
\newtcolorbox{boxNote}[1][]{title={{\emoji✨} ملاحظة\ifx\\#1\\\else ~#1\fi :}}
\newtcolorbox{boxAttention}[1][]{title={{\emoji🔔} تنبيه\ifx\\#1\\\else ~#1\fi :}}
\newtcolorbox{boxWarning}[1][]{title={{\emoji⚡} ملاحظة هامة\ifx\\#1\\\else ~#1\fi :}}
\newtcolorbox{boxSolution}[1][]{title={{\emoji✅} حل\ifx\\#1\\\else ~#1\fi :}}
\newtcolorbox{boxSymbol}[1][]{title={{\emoji🔣} رمز\ifx\\#1\\\else ~#1\fi :}}
\newtcolorbox{boxHint}[1][]{title={{\emoji💡} تلميح\ifx\\#1\\\else ~#1\fi :}}
\else
\tcbset{colback=white, colframe=black, fonttitle=\bfseries, boxrule=0.8pt}
\newtcolorbox{boxDef}[1][]{colback=blue!5!white,colframe=blue!75!black,
  title={{\emoji📘} تعريف\ifx\\#1\\\else ~#1\fi :}}
\newtcolorbox{boxExercise}[1][]{colback=cyan!5!white,colframe=cyan!70!black,
  title={{\emoji🧩} تمرين\ifx\\#1\\\else ~#1\fi :}}
\newtcolorbox{boxExample}[1][]{colback=yellow!5!white,colframe=orange!90!black,
  title={{\emoji📝} مثال\ifx\\#1\\\else ~#1\fi :}}
\newtcolorbox{boxNote}[1][]{colback=gray!10!white,colframe=black,
  title={{\emoji✨} ملاحظة\ifx\\#1\\\else ~#1\fi :}}
\newtcolorbox{boxAttention}[1][]{colback=magenta!10!white,colframe=magenta!80!black,
  title={{\emoji🔔} تنبيه\ifx\\#1\\\else ~#1\fi :}}
\newtcolorbox{boxWarning}[1][]{colback=red!5!white,colframe=red!75!black,
  title={{\emoji⚡} ملاحظة هامة\ifx\\#1\\\else ~#1\fi :}}
\newtcolorbox{boxSolution}[1][]{colback=green!5!white,colframe=green!60!black,
  title={{\emoji✅} حل\ifx\\#1\\\else ~#1\fi :}}
\newtcolorbox{boxSymbol}[1][]{colback=purple!5!white,colframe=purple!70!black,
  title={{\emoji🔣} رمز\ifx\\#1\\\else ~#1\fi :}}
\newtcolorbox{boxHint}[1][]{colback=teal!5!white,colframe=teal!60!black,
  title={{\emoji💡} تلميح\ifx\\#1\\\else ~#1\fi :}}
\fi


\ifwhiteblack
\tcbset{simplecode/.style={ colback=white, colframe=black, boxrule=1pt, arc=2pt, left=4pt,right=4pt,top=4pt,bottom=4pt}}
\else
\tcbset{simplecode/.style={ colback=gray!5, colframe=black!50, boxrule=0.4pt, arc=2pt, left=4pt,right=4pt,top=4pt,bottom=4pt}}
\fi
% Compact code boxes: tighter line spacing (the preamble uses onehalfspacing) and small vertical gaps.
% Must come AFTER the simplecode definition above.
\tcbset{simplecode/.append style={fontupper=\setstretch{0.95}, top=4pt, bottom=4pt, before skip=6pt, after skip=6pt}}
\newenvironment{boxCode}{\begin{tcolorbox}[simplecode]}{\end{tcolorbox}}

% Code cell used inside tables (keeps code left-to-right inside RTL tables).
% Use inside a fixed-width p{...} column: \begin{pycell}\begin{minted}...\end{minted}\end{pycell}
\newtcolorbox{pycell}{frame empty,colback=white,boxsep=0pt,left=2pt,right=2pt,top=2pt,bottom=2pt,width=\linewidth}

\newcolumntype{C}[1]{>{\centering\arraybackslash}p{#1}}

\makeatletter
\DTMnewdatestyle{dotted}{%
  \renewcommand{\DTMdisplaydate}[4]{\two@digits{##3}.\two@digits{##2}.##1}%
  \renewcommand{\DTMDisplaydate}{\DTMdisplaydate}%
}
\makeatother
\DTMsetdatestyle{dotted}

% redefine spaces after titles
\makeatletter
\renewcommand{\@maketitle}{%
  \begin{center}
    {\huge \bfseries \@title \par}%
    \vskip 0.2em % space between title and author
    {\large \@author \par}%
    % \vskip 0.2em % space between author and date
    % {\normalsize \@date \par}%
    {\small \hfill \begin{english}\DTMtoday\end{english} \par}%
  \end{center}
}
\makeatother

\fancyhf{} % clear default
\fancypagestyle{plain}{
  \fancyhf{}
  \fancyhead[L]{مدرسة التسامح الشاملة}
  % \fancyhead[L]{\includegraphics[height=1cm]{../../../images/logoTasamoh.png}}
  \fancyhead[R]{الأستاذ محمود اغبارية}
  \fancyfoot[C]{\thepage}
}

\fancyhead[L]{مدرسة التسامح الشاملة}
\fancyhead[R]{الأستاذ محمود اغبارية}
\fancyfoot[C]{\thepage}

\setcounter{tocdepth}{3} % only section subsection and subsubsection in TOC

\newcommand{\importpdfpage}[2]{%
    \thispagestyle{empty}
    \newgeometry{left=0cm,right=0cm,top=0cm,bottom=0cm}
    \includegraphics[page=#2,width=0.95\paperwidth,height=0.95\paperheight,keepaspectratio]{#1}
    \restoregeometry
}

\newcommand{\insertFullImg}[1]{%
    \noindent
    \makebox[\textwidth][c]{
        \includegraphics[width=0.9\paperwidth,keepaspectratio]{#1}%
    }%
}%


\ifwithsols
\title{حل ورقة تمرن 6 قسم 2 للصف العاشر $10$ \\ أسئلة تتبّع في جملة الشرط \textenglish{if}}
\else
\title{ورقة تمرن 6 قسم 2 للصف العاشر $10$ \\ أسئلة تتبّع في جملة الشرط \textenglish{if}}
\fi

\begin{document}

\maketitle
\thispagestyle{fancy}

\begin{enumerate}[itemsep=2.2em]

% ==============================================================
% السؤال 1: تتبّع بدون جملة شرط (متغيرات وحسابات وعمليات Math)
% ==============================================================
\item
أمامك مقطع البرنامج التالي الذي يستقبل عددًا صحيحًا موجبًا $num$ مكوّنًا من $3$ منازل (بين $100$ و $999$):

\begin{boxCode}
\begin{english}
\begin{minted}{csharp}
int num = int.Parse(Console.ReadLine());
int a = num / 100;
int b = (num / 10) % 10;
int c = num % 10;
int m = Math.Max(a, c);
int d = Math.Abs(a - b);
int res = m * 10 + d;
Console.WriteLine(res);
\end{minted}
\end{english}
\end{boxCode}

\begin{enumerate}[itemsep=1.8em]
    \filbreak
    \item \textbf{ما هو المُخرج؟} \\
    تتبّع تنفيذ المقطع البرمجي أعلاه لكل قيمة من قيم المُدخل $num$ في الجدول التالي، واكتب قيم المتغيرات والمُخرج الذي يطبعه البرنامج:
    \begin{center}
    \begin{english}
    \renewcommand{\arraystretch}{1.45}
    \begin{tabular}{|C{1.8cm}|C{1.2cm}|C{1.2cm}|C{1.2cm}|C{1.2cm}|C{1.2cm}|C{2.2cm}|}
    \hline
    \textbf{num} & \textbf{a} & \textbf{b} & \textbf{c} & \textbf{m} & \textbf{d} & \textarabic{\textbf{المُخرج}} \\
    \hline
    472 & & & & & & \\
    \hline
    389 & & & & & & \\
    \hline
    661 & & & & & & \\
    \hline
    515 & & & & & & \\
    \hline
    \end{tabular}
    \end{english}
    \end{center}
    \filbreak

    \item \textbf{أعطِ مُدخلاً يطبع كلّ واحد من هذه المخرجات:} \\
    لكل مُخرج في الجدول التالي، اقترح عددًا صحيحًا $num$ مكوّنًا من $3$ منازل (\textbf{مختلفًا} عن أعداد الفرع أ) يجعل البرنامج يطبع هذا المُخرج، مع تحقيق الشرط الإضافي المذكور:
    \begin{center}
    \renewcommand{\arraystretch}{1.45}
    \begin{tabular}{|C{2.4cm}|p{7cm}|C{2.4cm}|}
    \hline
    \textbf{المُخرج المطلوب} & \textbf{شرط إضافي على المُدخلات} & \textbf{قيمة $num$} \\
    \hline
    $83$ & بحيث تكون منزلة المئات أكبر من منزلة الآحاد ($a > c$) & \\
    \hline
    $83$ & بحيث تكون منزلة المئات أصغر من منزلة الآحاد ($a < c$) & \\
    \hline
    $60$ & بحيث تكون جميع منازل العدد غير صفرية & \\
    \hline
    $75$ & بحيث تكون منزلة المئات أصغر من منزلة العشرات ($a < b$) & \\
    \hline
    \end{tabular}
    \end{center}
    \filbreak
\end{enumerate}

\ifwithsols
\begin{boxSolution}[الفرع أ والفرع ب]
\textbf{حل الفرع أ:}
\begin{center}
\begin{english}
\renewcommand{\arraystretch}{1.35}
\begin{tabular}{|C{1.8cm}|C{1.2cm}|C{1.2cm}|C{1.2cm}|C{1.2cm}|C{1.2cm}|C{2.2cm}|}
\hline
\textbf{num} & \textbf{a} & \textbf{b} & \textbf{c} & \textbf{m} & \textbf{d} & \textarabic{\textbf{المُخرج}} \\
\hline
472 & 4 & 7 & 2 & 4 & 3 & 43 \\
\hline
389 & 3 & 8 & 9 & 9 & 5 & 95 \\
\hline
661 & 6 & 6 & 1 & 6 & 0 & 60 \\
\hline
515 & 5 & 1 & 5 & 5 & 4 & 54 \\
\hline
\end{tabular}
\end{english}
\end{center}

\textbf{حل الفرع ب (أمثلة لإجابات صحيحة):}
\begin{itemize}[nosep]
    \item لطباعة $83$ مع $a > c$: يجب أن يكون $m = 8$ و $d = |a - b| = 3$، وبما أن $a > c$ فإن $a = 8$ و $b = 5$ و $c < 8$، مثل: $num = 852$.
    \item لطباعة $83$ مع $a < c$: يجب أن يكون $c = 8$ و $a < 8$ و $|a - b| = 3$، مثل: $num = 478$ (أو $528$).
    \item لطباعة $60$ بمنازل غير صفرية: يجب أن يكون $m = 6$ و $d = 0$ (أي $a = b$)، مثل: $num = 664$ (أو $336$).
    \item لطباعة $75$ مع $a < b$: يجب أن يكون $m = 7$ و $b - a = 5$، مثل: $num = 277$ (أو $387$ أو $497$).
\end{itemize}
\end{boxSolution}
\fi
\filbreak

% ==============================================================
% السؤال 2: تتبّع مع Math وجملة if - else وشروط مركّبة
% ==============================================================
\item
أمامك مقطع البرنامج التالي الذي يستقبل عددين صحيحين $a$ و $b$:

\begin{boxCode}
\begin{english}
\begin{minted}{csharp}
int a = int.Parse(Console.ReadLine());
int b = int.Parse(Console.ReadLine());
int max = Math.Max(a, b);
int sum = a + b;

if (max >= 10 && (sum % 2 == 0 || a < 0))
{
    Console.WriteLine("A");
}
else
{
    Console.WriteLine("B");
}
\end{minted}
\end{english}
\end{boxCode}

\begin{enumerate}[itemsep=1.8em]
    \filbreak
    \item \textbf{ما هو المُخرج؟} \\
    تتبّع تنفيذ المقطع البرمجي أعلاه لكل زوج من المدخلات في الجدول التالي، واكتب قيمتَي المتغيرين $max$ و $sum$ والمُخرج الذي يطبعه البرنامج:
    \begin{center}
    \begin{english}
    \renewcommand{\arraystretch}{1.45}
    \begin{tabular}{|C{2cm}|C{2cm}|C{2cm}|C{2cm}|C{2.8cm}|}
    \hline
    \textbf{a} & \textbf{b} & \textbf{max} & \textbf{sum} & \textarabic{\textbf{المُخرج}} \\
    \hline
    12 & 4 & & & \\
    \hline
    -5 & 10 & & & \\
    \hline
    10 & 3 & & & \\
    \hline
    -4 & 8 & & & \\
    \hline
    -12 & -3 & & & \\
    \hline
    \end{tabular}
    \end{english}
    \end{center}
    \filbreak

    \item \textbf{أعطِ مُدخلاً يطبع كلّ واحد من هذه المخرجات:} \\
    لكل مُخرج في الجدول التالي، اقترح زوجًا من الأعداد الصحيحة $a$ و $b$ (\textbf{مختلفًا} عن أزواج الفرع أ) يجعل البرنامج يطبع هذا المُخرج، مع تحقيق الشرط الإضافي المذكور:
    \begin{center}
    \renewcommand{\arraystretch}{1.45}
    \begin{tabular}{|C{2.4cm}|p{5.5cm}|C{1.9cm}|C{1.9cm}|}
    \hline
    \textbf{المُخرج المطلوب} & \textbf{شرط إضافي على المُدخلات} & \textbf{قيمة $a$} & \textbf{قيمة $b$} \\
    \hline
    $A$ & بحيث يكون العددان موجبَين ($a > 0$ و $b > 0$) & & \\
    \hline
    $A$ & بحيث يكون $sum$ عددًا فرديًا & & \\
    \hline
    $B$ & بحيث يكون $max \ge 10$ & & \\
    \hline
    $B$ & بحيث يكون $a < 0$ ويكون $sum$ عددًا زوجيًا & & \\
    \hline
    \end{tabular}
    \end{center}
    \filbreak
\end{enumerate}

\ifwithsols
\begin{boxSolution}[الفرع أ والفرع ب]
\textbf{حل الفرع أ:}
\begin{center}
\begin{english}
\renewcommand{\arraystretch}{1.35}
\begin{tabular}{|C{2cm}|C{2cm}|C{2cm}|C{2cm}|C{2.8cm}|}
\hline
\textbf{a} & \textbf{b} & \textbf{max} & \textbf{sum} & \textarabic{\textbf{المُخرج}} \\
\hline
12 & 4 & 12 & 16 & A \\
\hline
-5 & 10 & 10 & 5 & A \\
\hline
10 & 3 & 10 & 13 & B \\
\hline
-4 & 8 & 8 & 4 & B \\
\hline
-12 & -3 & -3 & -15 & B \\
\hline
\end{tabular}
\end{english}
\end{center}

\textbf{حل الفرع ب (أمثلة لإجابات صحيحة):}
\begin{itemize}[nosep]
    \item لطباعة $A$ بعددين موجبين: يجب أن يكون $max \ge 10$ والمجموع زوجيًا، مثل: $a = 10,\ b = 6$.
    \item لطباعة $A$ مع كون $sum$ فرديًا: يجب أن يكون $a < 0$ و $b \ge 10$ ومجموعهما فرديًا، مثل: $a = -3,\ b = 12$.
    \item لطباعة $B$ مع $max \ge 10$: يجب أن يكون $a \ge 0$ والمجموع فرديًا، مثل: $a = 2,\ b = 11$.
    \item لطباعة $B$ مع $a < 0$ و $sum$ زوجي: يجب إخلال الشرط $max \ge 10$ (أي $max < 10$)، مثل: $a = -2,\ b = 6$.
\end{itemize}
\end{boxSolution}
\fi
\filbreak

% ==============================================================
% السؤال 3: تتبّع مع Math وجملة if - else if - else وشروط مركّبة
% ==============================================================
\item
أمامك مقطع البرنامج التالي الذي يستقبل عددين صحيحين $x$ و $y$:

\begin{boxCode}
\begin{english}
\begin{minted}{csharp}
int x = int.Parse(Console.ReadLine());
int y = int.Parse(Console.ReadLine());
int p = (int)Math.Pow(x, 2) - y;

if (p < 0)
{
    Console.WriteLine("A");
}
else if (p == 0 || x == y)
{
    Console.WriteLine("B");
}
else if (!(x > 0 && y > 0))
{
    Console.WriteLine("C");
}
else
{
    Console.WriteLine("D");
}
\end{minted}
\end{english}
\end{boxCode}

\begin{enumerate}[itemsep=1.8em]
    \filbreak
    \item \textbf{ما هو المُخرج؟} \\
    تتبّع تنفيذ المقطع البرمجي أعلاه لكل زوج من المدخلات في الجدول التالي، واكتب قيمة المتغير $p$ والمُخرج الذي يطبعه البرنامج:
    \begin{center}
    \begin{english}
    \renewcommand{\arraystretch}{1.45}
    \begin{tabular}{|C{2.4cm}|C{2.4cm}|C{2.4cm}|C{3.2cm}|}
    \hline
    \textbf{x} & \textbf{y} & \textbf{p} & \textarabic{\textbf{المُخرج}} \\
    \hline
    2 & 7 & & \\
    \hline
    3 & 9 & & \\
    \hline
    4 & 4 & & \\
    \hline
    -2 & 1 & & \\
    \hline
    3 & 5 & & \\
    \hline
    \end{tabular}
    \end{english}
    \end{center}
    \filbreak

    \item \textbf{أعطِ مُدخلاً يطبع كلّ واحد من هذه المخرجات:} \\
    لكل مُخرج في الجدول التالي، اقترح زوجًا من الأعداد الصحيحة $x$ و $y$ (\textbf{مختلفًا} عن أزواج الفرع أ) يجعل البرنامج يطبع هذا المُخرج، مع تحقيق الشرط الإضافي المذكور:
    \begin{center}
    \renewcommand{\arraystretch}{1.45}
    \begin{tabular}{|C{2.4cm}|p{5.5cm}|C{1.9cm}|C{1.9cm}|}
    \hline
    \textbf{المُخرج المطلوب} & \textbf{شرط إضافي على المُدخلات} & \textbf{قيمة $x$} & \textbf{قيمة $y$} \\
    \hline
    $A$ & بحيث يكون $x < 0$ & & \\
    \hline
    $B$ & بحيث يكون $x \ne y$ & & \\
    \hline
    $B$ & بحيث يكون $p > 0$ & & \\
    \hline
    $C$ & بحيث يكون $x > 0$ & & \\
    \hline
    $D$ & بحيث يكون $x < y$ & & \\
    \hline
    \end{tabular}
    \end{center}
    \filbreak
\end{enumerate}

\ifwithsols
\begin{boxSolution}[الفرع أ والفرع ب]
\textbf{حل الفرع أ:}
\begin{center}
\begin{english}
\renewcommand{\arraystretch}{1.35}
\begin{tabular}{|C{2.4cm}|C{2.4cm}|C{2.4cm}|C{3.2cm}|}
\hline
\textbf{x} & \textbf{y} & \textbf{p} & \textarabic{\textbf{المُخرج}} \\
\hline
2 & 7 & -3 & A \\
\hline
3 & 9 & 0 & B \\
\hline
4 & 4 & 12 & B \\
\hline
-2 & 1 & 3 & C \\
\hline
3 & 5 & 4 & D \\
\hline
\end{tabular}
\end{english}
\end{center}

\textbf{حل الفرع ب (أمثلة لإجابات صحيحة):}
\begin{itemize}[nosep]
    \item لطباعة $A$ مع $x < 0$: يجب أن يكون $x^2 < y$، مثل: $x = -2,\ y = 10$ (حيث $p = -6$).
    \item لطباعة $B$ مع $x \ne y$: يجب أن يكون $p = 0$ (أي $y = x^2$)، مثل: $x = 2,\ y = 4$ (أو $x = 5,\ y = 25$).
    \item لطباعة $B$ مع $p > 0$: يجب أن يكون $x = y$ و $x^2 - x > 0$، مثل: $x = 3,\ y = 3$ (حيث $p = 6$).
    \item لطباعة $C$ مع $x > 0$: يجب أن يكون $y \le 0$، مثل: $x = 2,\ y = -3$.
    \item لطباعة $D$ مع $x < y$: يجب أن يكون $x > 0$ و $y > 0$ و $x < y < x^2$، مثل: $x = 3,\ y = 6$ (حيث $p = 3$).
\end{itemize}
\end{boxSolution}
\fi
\filbreak

% ==============================================================
% السؤال 4: تتبّع مع Math وجمل شرط متداخلة (Nested if)
% ==============================================================
\item
أمامك مقطع البرنامج التالي الذي يستقبل عددين صحيحين $a$ و $b$:

\begin{boxCode}
\begin{english}
\begin{minted}{csharp}
int a = int.Parse(Console.ReadLine());
int b = int.Parse(Console.ReadLine());
int min = Math.Min(a, b);
int max = Math.Max(a, b);

if (min <= 0)
{
    if (max > 0 && Math.Abs(min) > max)
    {
        Console.WriteLine("A");
    }
    else
    {
        Console.WriteLine("B");
    }
}
else
{
    if (max % min == 0 && a != b)
    {
        Console.WriteLine("C");
    }
    else
    {
        Console.WriteLine("D");
    }
}
\end{minted}
\end{english}
\end{boxCode}

\begin{enumerate}[itemsep=1.8em]
    \filbreak
    \item \textbf{ما هو المُخرج؟} \\
    تتبّع تنفيذ المقطع البرمجي أعلاه لكل زوج من المدخلات في الجدول التالي، واكتب قيمتَي المتغيرين $min$ و $max$ والمُخرج الذي يطبعه البرنامج:
    \begin{center}
    \begin{english}
    \renewcommand{\arraystretch}{1.45}
    \begin{tabular}{|C{2cm}|C{2cm}|C{2cm}|C{2cm}|C{2.8cm}|}
    \hline
    \textbf{a} & \textbf{b} & \textbf{min} & \textbf{max} & \textarabic{\textbf{المُخرج}} \\
    \hline
    -8 & 5 & & & \\
    \hline
    -3 & 7 & & & \\
    \hline
    0 & -4 & & & \\
    \hline
    4 & 12 & & & \\
    \hline
    6 & 6 & & & \\
    \hline
    \end{tabular}
    \end{english}
    \end{center}
    \filbreak

    \item \textbf{أعطِ مُدخلاً يطبع كلّ واحد من هذه المخرجات:} \\
    لكل مُخرج في الجدول التالي، اقترح زوجًا من الأعداد الصحيحة $a$ و $b$ (\textbf{مختلفًا} عن أزواج الفرع أ) يجعل البرنامج يطبع هذا المُخرج، مع تحقيق الشرط الإضافي المذكور:
    \begin{center}
    \renewcommand{\arraystretch}{1.45}
    \begin{tabular}{|C{2.4cm}|p{5.5cm}|C{1.9cm}|C{1.9cm}|}
    \hline
    \textbf{المُخرج المطلوب} & \textbf{شرط إضافي على المُدخلات} & \textbf{قيمة $a$} & \textbf{قيمة $b$} \\
    \hline
    $A$ & بحيث يكون $a > 0$ & & \\
    \hline
    $B$ & بحيث يكون $max > 0$ وكلا العددين غير صفري & & \\
    \hline
    $C$ & بحيث يكون $a > b$ & & \\
    \hline
    $D$ & بحيث يقبل $max$ القسمة على $min$ بدون باقٍ & & \\
    \hline
    $D$ & بحيث يكون $a \ne b$ & & \\
    \hline
    \end{tabular}
    \end{center}
    \filbreak
\end{enumerate}

\ifwithsols
\begin{boxSolution}[الفرع أ والفرع ب]
\textbf{حل الفرع أ:}
\begin{center}
\begin{english}
\renewcommand{\arraystretch}{1.35}
\begin{tabular}{|C{2cm}|C{2cm}|C{2cm}|C{2cm}|C{2.8cm}|}
\hline
\textbf{a} & \textbf{b} & \textbf{min} & \textbf{max} & \textarabic{\textbf{المُخرج}} \\
\hline
-8 & 5 & -8 & 5 & A \\
\hline
-3 & 7 & -3 & 7 & B \\
\hline
0 & -4 & -4 & 0 & B \\
\hline
4 & 12 & 4 & 12 & C \\
\hline
6 & 6 & 6 & 6 & D \\
\hline
\end{tabular}
\end{english}
\end{center}

\textbf{حل الفرع ب (أمثلة لإجابات صحيحة):}
\begin{itemize}[nosep]
    \item لطباعة $A$ مع $a > 0$: يجب أن يكون $b < 0$ و $|b| > a$، مثل: $a = 4,\ b = -7$.
    \item لطباعة $B$ مع $max > 0$ وعددين غير صفريين: يجب أن يكون أحدهما سالبًا والآخر موجبًا مع $|min| \le max$، مثل: $a = -2,\ b = 5$.
    \item لطباعة $C$ مع $a > b$: يجب أن يكون كلاهما موجبًا و $a$ من مضاعفات $b$، مثل: $a = 15,\ b = 5$.
    \item لطباعة $D$ مع قابلية قسمة $max$ على $min$: يجب أن يكون العددان موجبين ومتساويين ($a = b$)، مثل: $a = 5,\ b = 5$.
    \item لطباعة $D$ مع $a \ne b$: يجب أن يكون كلاهما موجبًا ولا يقبل الأكبر القسمة على الأصغر، مثل: $a = 7,\ b = 4$.
\end{itemize}
\end{boxSolution}
\fi
\filbreak

\end{enumerate}

\vspace{1cm}
\begin{flushleft}
أرجو لكم وقتًا ممتعًا.

الأستاذ محمود اغبارية.
\end{flushleft}

\end{document}
